% !TEX root = ../main.tex

\documentclass[../main.tex]{subfiles}

\graphicspath{{\subfix{../images/}}}

\begin{document}

\subsection{Установка зависимостей}
 
\begin{lstlisting}[language=bash,caption={Установка системных зависимостей}]
sudo apt update
sudo apt install build-essential cmake git \
    libglfw3-dev libglew-dev libglm-dev libgl1-mesa-dev
\end{lstlisting}
 
Библиотека MAVLink является заголовочной и клонируется напрямую в каталог
проекта:
 
\begin{lstlisting}[language=bash,caption={Клонирование MAVLink}]
cd pixhawk_imu_visualizer
git clone https://github.com/mavlink/c_library_v2.git
\end{lstlisting}
 
\subsection{Сборка проекта}
 
\begin{lstlisting}[language=bash,caption={Сборка через CMake}]
mkdir build && cd build
cmake ..
make -j$(nproc)
\end{lstlisting}
 
При успешной сборке в каталоге \texttt{build/} появляется исполняемый файл
\texttt{pixhawk\_imu\_reader}.
 
\subsection{Настройка прав и запуск}
 
\begin{lstlisting}[language=bash,caption={Права на порт и запуск}]
sudo usermod -a -G dialout $USER

sudo chmod 666 /dev/ttyACM0
 
cd build
./pixhawk_imu_reader
\end{lstlisting}
 
Окно визуализации открывается немедленно. Для завершения~--- Ctrl+C в
терминале.
 
\subsection{Управление в окне визуализации}
 
\begin{table}[H]
\centering
\caption{Управление в окне визуализации}
\label{tab:controls}
\begin{tabularx}{\textwidth}{|l|X|}
\hline
\textbf{Действие} & \textbf{Эффект} \\
\hline
ЛКМ + перетаскивание & Вращение камеры вокруг объекта \\
\hline
Колесо мыши          & Приближение / отдаление (0{,}5~---~30~м) \\
\hline
Клавиша \textbf{R}   & Сброс позиции и траектории в $(0,\,0,\,0)$ \\
\hline
\textbf{Ctrl+C}      & Завершение программы \\
\hline
\end{tabularx}
\end{table}
 
\newpage
	
\end{document}
